\section{Toom's Contours}

\subsection{The Barrier Argument}


We are going to move between the infinite and the finite back and forth, so we will need to extend the size function, and the reduced notion to chains on the finite-complex.  
\begin{definition}
Consider a $d^{\prime}$-dimensional chain on the $d$-dimensional Toric $x \in (\mathcal{F}^{n}_{d^{\prime}}; \mathbb{F}_{2})$. We extend the notation of size to $x$ as follows. Let $x^{\prime}$ be the $\infty$-lift of $x$, so $x^{\prime} \in (\mathcal{F}_{d^{\prime}}; \mathbb{F}_{2})$. Let the chain $y \in (\mathcal{F}_{d^{\prime}}; \mathbb{F}_{2})$ be a subset chain of $x^{\prime}$ such:
  \begin{enumerate}
    \item The projection of $y$ back to finite complex covers $x$.
    \item The size of $y$ is minimal. 
  \end{enumerate}
  Then, we set: 
  \begin{equation*}
    \begin{split}
      \Size{x} \colonequals \Size{y}
    \end{split}
  \end{equation*}
\end{definition}


\begin{claim}
  \inb{\textbf{Triangle ineq for finite chains. Prove it. } }
\end{claim}

\begin{claim}
Suppose that the chain $x \in (\mathcal{F}^{n}_{d^{\prime}}; \mathbb{F}_{2})$ closes a non-trivial loop on the Toric, then $\Size{x} \ge n$.
\end{claim}
\begin{proof}
Let $x^{\prime}$ be the $\infty$-lift of $x$, and consider a subset $y \subset x^{\prime}$ such that the projection of $y$ covers $x$. In particular, there are vertices $v_{1}, \ldots, v_{l}$ in $y$ such that their projection corresponds to the non-trivial loop in $x$. Assume the non-trivial cycle is in the $g$-direction, so the $L_{g}$ values of the projections of $v_{1}$ and $v_{l}$ are $0$ and $n-1$, respectively. Without loss of generality, assume that $L_{g}(v_{l}) > L_{g}(v_{1})$. Thus:
  \begin{equation*}
    \begin{split}
      \Size{y} \ge \Size{\{v_{1},v_{l}\}} \ge \sum_{g^{\prime}\neq g} L_{g^{\prime}}(v_{1}) + L_{g}(v_{l}) + L_{d+1}(v_{1}) = L_{g}(v_{l}) -L_{g}(v_{1})
    \end{split}
  \end{equation*}
And it's trivial to see that the assumption $L_{g}(v_{l}) > L_{g}(v_{1})$ implies that the difference is at least $n-1$.
\end{proof}

\begin{claim}
  \label{claim:fillchain}
  Let $x \in (\mathcal{F}^{n}_{d^{\prime}}; \mathbb{F}_{2})$ be a $d^{\prime}$-dimensional chain, such that the $d^{\prime}-1$-dimensional chain $\partial^{-} x$ is not connected. Suppose that for some constant number $\alpha \in (0,1)$, each connected component of $\partial^{-}x$ has size at most $\alpha n$. Then there are $d^{\prime}$-dimensional chains $y_{1}, y_{2}, \ldots$ such that each $\partial^{-}y_{i}$ is a connected component of $\partial^{-}x$ and $\Size{y_{i}} \le \alpha n$.

\end{claim}

\begin{proof}
  Consider a $d^{\prime}$-dimensional chain $x$ such that $\partial^{-} x$ has a decomposition $x = \tilde{x}_{1} + \tilde{x}_{2}$, where $\tilde{x}_{1}, \tilde{x}_{2}$ are $d^{\prime}-1$-dimensional chains, and $\tilde{x}_{1}$ is a connected component of $\partial^{-}x$. In particular $\tilde{x}_{1}$ is not connected to $\tilde{x}_{2}$. By disconnectivity, $\tilde{x}_{1}$ does not share a $d^{\prime}-2$ dimensional face with $\tilde{x}_{2}$. Thus, $\partial^{-} \tilde{x}_{1}$ and $\partial^{-} \tilde{x}_{2}$ are disjoint. On the other hand, since $\partial^{- \circ 2} x = 0$, then combining with disjointness, it follows that:
  \begin{equation*}
    \begin{split}
      \partial^{-}( \tilde{x}_{1} + \tilde{x}_{2} ) = 0 \Rightarrow \partial^{-} \tilde{x}_{1} = \partial^{-}\tilde{x}_{2} = 0 
    \end{split}
  \end{equation*}  Namely, $\tilde{x}_{1} \in \ker \partial^{-} : (\mathcal{F}^{n}_{d^{\prime}-1}; \mathbb{F}_{2}) \rightarrow ( \mathcal{F}^{n}_{d^{\prime}-2}; \mathbb{F}_{2})$.
Suppose $\tilde{x}_{1} \notin \image \partial^{-} : ( \mathcal{F}^{n}_{d^{\prime}} ; \mathbb{F}_{2}) \rightarrow ( \mathcal{F}^{n}_{d^{\prime}-1}; \mathbb{F}_{2})$, then we have that $\tilde{x}_{1}$ is a non-trivial codeword of the $(d,d^{\prime}-1)$-Toric code. 
Since $d^{\prime} \ge 2$, the size of $\tilde{x}$ is at least $n$, which is a contradiction. So we get that $\tilde{x}_{1}\in \image \partial^{-} : ( \mathcal{F}^{n}_{d^{\prime}} ; \mathbb{F}_{2}) \rightarrow ( \mathcal{F}^{n}_{d^{\prime}-1}; \mathbb{F}_{2})$. In particular, there is a $d^{\prime}$-dimensional chain $x_{1}$ such that $\partial^{-}x_{1} = \tilde{x}_{1}$. Choose $x_{1}$ such it doesn't supported on a non-trivial codeword, namely $\tilde{x}_{1}$ is finite. The reduced form of $x_{1}$ has size at most $\Size{\tilde{x}_{1}} \le \alpha n$, and satisfies the desired condition for the connected component $\tilde{x}_{1}$. To complete, we repeat recursively on the connected components of $\tilde{x}_{2}$.
\end{proof}


\begin{definition}
Let $x \in (\mathcal{F}^{n}_{d^{\prime}}; \mathbb{F}_{2})$ be a $d^{\prime}$-dimensional chain. $x$ will be said to be a fillable chain if each of the connected components of its $\partial^{-}$-boundary, namely $\partial^{-} x$, has size at most $n$. By \Cref{claim:fillchain}, there are $d^{\prime}$-dimensional reduced chains $x_{1}, \ldots,$ such that each $\partial^{-}x_{i}$ is a connected component of $\partial^{-}x$. Thus:
    \begin{equation*}
      \begin{split}
      \partial^{-} (x + \sum_{i}x_{i}) = 0  \quad \Rightarrow x + \sum_{i}x_{i} \in \ker \partial^{-} \colon (\mathcal{F}^{n}_{d^{\prime}}; \mathbb{F}_{2}) \rightarrow ( \mathcal{F}^{n}_{d^{\prime}-1}; \mathbb{F}_{2})
    \end{split}
  \end{equation*}
Namely, $x + \sum_{i}x_{i}$ is a codeword in $\mathcal{C}_{X}/ \mathcal{C}_{Z}^{\perp}$. We define the filling operator, denoted by $[x]$, which for a given fillable chain $x$, returns the homology class (logical operator) of $x + \sum_{i}x_{i}$.
\end{definition}


%\begin{claim}
%  Let $x_{1},x_{2} \in (\mathcal{F}^{n}_{d^{\prime}}; \mathbb{F}_{2})$ be $d^{\prime}$-dimensional chains, such the Hamming distance between $x_{1}$ and $x_{2}$ is one. Then: 
%  \begin{equation*}
%    \begin{split}
%      \Size{[x_{1}] + [x_{2}]} 
%    \end{split}
%  \end{equation*}
%\end{claim}

\begin{claim}
  \label{claim:mergebyf}
  Let $x_{1},x_{2} \in (\mathcal{F}^{n}_{d^{\prime}}; \mathbb{F}_{2})$ be $d^{\prime}$-dimensional chains, and let $l$ be an integer such the Hamming distance between $x_{1}$ and $x_{2}$ is at most $l$. Denote by $r$ the size of the maximal connected component of $\partial^{-}x_{1}$. Then the maximal size of a connected component of $\partial^{-}x_{2}$ is at most ${ d^{\prime} \choose 2 }\Delta r$
\end{claim}
\begin{proof}
  We argue that each connected component of $\partial^{-}x^{\prime}_{i}$ has a size of at most ${ d^{\prime} \choose 2 }\Delta c n$. To see this, consider the adjacency graph (\cref{def:asontoric}) respecting $d^{\prime}-1$-dimensional chains, where each edge in the graph is a $d^{\prime}-2$ chain. Introducing a new $d^{\prime}$-dimensional face adds at most ${ d^{\prime} \choose 2 }$ edges, each corresponding to a $d^{\prime}-2$ face contained in the new $d^{\prime}$-dimensional face. So introducing $\Delta$ faces into $x^{\prime}_{i}$ adds at most ${ d^{\prime} \choose 2 } \Delta$ new edges, and therefore connects at most the same number of connected components in $\partial^{-}x^{\prime}_{i}$. Using the triangle inequality for intersecting sets (\Cref{claim:tringinq}), we have that any connected component of $\partial^{-}x_{j + i\Delta}^{\prime}$ is at most ${ d^{\prime} \choose 2 } \Delta$ times the maximal size of a connected component of $\partial^{-}x_{1 + i\Delta}^{\prime}$, namely ${ d^{\prime} \choose 2 } \Delta c n$.
\end{proof}

\begin{claim}
  \label{claim:barrier}
Let $x_{1}, \ldots, x_{t} \in (\mathcal{F}^{n}_{d^{\prime}}; \mathbb{F}_{2})$ be $d^{\prime}$-dimensional chains, such that $x_{1}$ is a trivial chain and for any $i \in [t-1]$, the Hamming distance between $x_{i}$ and $x_{i+1}$ is one. There is a constant $c \in (0,1)$, depending only on $d, d^{\prime}$, such that if $n$ is sufficiently large, and for any $x_{i} \in x_{1}, \ldots, x_{t}$, the size of each of the connected components of $\partial^{-} x_{i}$ is bounded by $c n$, then:
  \begin{equation*}
    \begin{split}
      [0]= [x_{1}] = [x_{2}] = ... = [x_{t}] 
    \end{split}
  \end{equation*}
\end{claim}
\begin{proof}
Assume through contradiction that there is $i$ such that $0 = [x_{i}] \neq [x_{i+1}]$. Denote by $x_{i}^{1}, x_{i}^{2}, \ldots$ and $x_{i+1}^{1}, x_{i+2}^{2}, \ldots$ the reduced $d^{\prime}$-dimensional chains such that $\partial^{-}x_{i}^{j}$ is a connected component of $\partial^{-}x_{i}$. Hence:
  \begin{equation*}
    \begin{split}
      [x_{i+1}] \sim_{\mathcal{C}_{Z^{\perp}}} [x_{i+1}] + [x_{i}] \sim_{\mathcal{C}_{Z^{\perp}}}  x_{i} + x_{i+1} + \sum_{j}x_{i}^{j} + \sum_{j}x_{i+1}^{j+1}
    \end{split}
  \end{equation*}
Denote by $w$ the chain:
  \begin{equation*}
    \begin{split}
      w \colonequals x_{i} + x_{i+1} + \sum_{j}x_{i}^{j} + \sum_{j}x_{i+1}^{j+1}
    \end{split}
  \end{equation*}
  Since the Hamming distance between $x_{i}$ and $x_{i+1}$ is at most one, and since each $d^{\prime}$-dimensional face contributes a constant support over $d^{\prime}-1$-dimensional chains, the $\partial^{-}$-boundaries of $x_{i}$ and $x_{i+1}$ are at a Hamming distance that is at most constant. Therefore, there are only constant connected components of $\partial^{-}x_{i}$ which are not connected components in $\partial^{-}x_{i+1}$ and vice versa. Denote by $y_{1}, \ldots, y_{m}$ the chains in $x_{i}^{1}, \ldots$ and $x_{i+1}^{1}, \ldots$ such that they belong to only one of the collections. Where $m$ is bounded by a function of $d,d^{\prime}$, namely $m \le M(d,d^{\prime})$. So:
  \begin{equation*}
    \begin{split}
      \sum_{j}x_{i}^{j} + \sum_{j}x_{i+1}^{j+1} = \sum_{j}^{m}y_{j}
    \end{split}
  \end{equation*}
\inb{\textbf{Do something with that line. } }  Moreover, the vertices of $y_{1}, \ldots, y_{m}$ and $x_{i+1}+x_{i}$ have a nontrivial intersection. This is because the face $x_{i+1}+x_{i}$ is responsible for the splits and merges of the $x_{i}$'s $\partial^{-}$-boundaries. Therefore:
\begin{equation*}
  \begin{split}
    \Size{ w } \le \Size{  x_{i} + x_{i+1} } + \sum_{j}^{m} \Size{y_{j}} \le d + M(d,d^{\prime}) c n  
  \end{split}
\end{equation*}
Thus, for sufficiently large $n$ ,if $c < \frac{1}{M(d,d^{\prime})+1}$, we have that $\Size{w} < n$, a contradiction for $w$ being a non-trivial codeword.
\end{proof}



\begin{claim}
Let $x_{1}, \ldots, x_{t} \in (\mathcal{F}^{n}_{d^{\prime}}; \mathbb{F}_{2})$ be $d^{\prime}$-dimensional chains defined as in \Cref{claim:barrier}, with the change that any two sequential chains $x_{i}, x_{i+1}$ are at a Hamming distance of at most $\Delta$, and setting the constant $c \leftarrow c/\Delta$. Then the claim still holds.
\end{claim}

\begin{proof}
  Let $x_{1}^{\prime}, \ldots, x^{\prime}_{\Delta t + 1}$ be the chains obtained from $x_{1}, \ldots, x_{t}$ by taking the stacking between any two sequential elements $\Delta-1$ chains, such that in the new set any two sequential chains are at Hamming distance at most one. Namely, for any $i \in [t]$, we have that $x_{1 + i\Delta}^{\prime} = x_{i + i}$ and $x^{\prime}_{i+1}$ differs from $x^{\prime}_{i}$ by at most one coordinate. By \Cref{claim:mergebyf}, each connected component of $\partial^{-}x^{\prime}_{i}$ has a size of at most ${ d^{\prime} \choose 2 }\Delta c n$. 

  \paragraph{}
  Thus, by \Cref{claim:barrier},we have that:
  \begin{equation*}
    \begin{split}
      [0] = [x_{1}^{\prime}] =[x_{2}^{\prime}] = \ldots = [x_{\Delta t + 1}^{\prime}] 
    \end{split}
  \end{equation*}
  In particular: 
  \begin{equation*}
    \begin{split}
      & [0] = [x_{1}^{\prime}] =[x_{1 + \Delta}^{\prime}] = \ldots  =  [x_{1 + i \Delta}^{\prime} ]  =  \ldots = [x_{\Delta t + 1}^{\prime}] \\
      \Rightarrow  & [0] = [x_{1}] =[x_{2}] =   \ldots = [x_{t}] \\
    \end{split}
  \end{equation*}
\end{proof}

\begin{corollary}[Barrier argument]
  \label{coro:barrier}
  Barrier argument.
\end{corollary}
